\section{Algorithms}
\label{sec:algos}

Both algorithms share the same outer loop, the same rollout mechanics, and the
same reward. For each of the $B=15$ sources in a batch, a fresh
$\lmb \sim \mathcal{U}[0,1]$ is drawn and $\Ggrp = 16$ prompts are sampled
i.i.d.\ from $\pol(\cdot \mid \lmb)$, forming a \emph{group} that shares a
source and a $\lmb$. Each prompt's reward $r_i = \hat{r}_{\lmb}(z_i, x)$ is the
Monte Carlo average of Eq.~\eqref{eq:mcreward}. The algorithms differ only in
how a group of $(z_i, r_i)$ pairs becomes a loss. Throughout, the policy
log-probability of a prompt is the token sum
\begin{equation}
  \log \pol(z \mid \lmb) \;=\; \sum_{t=1}^{L} \log \pol(z_t \mid z_{<t}, \lmb),
\end{equation}
and one crucial implementation fact holds for both methods: \textbf{exactly one
gradient update is taken per rollout}. There is no inner epoch loop.

\subsection{Baseline: GRPO}
\label{sec:grpo}

GRPO replaces a learned critic with a group-relative advantage. Within each
group, rewards are standardized,
\begin{equation}
  A_i \;=\; \frac{r_i - \mathrm{mean}(r_{1:\Ggrp})}
                 {\mathrm{std}(r_{1:\Ggrp}) + \varepsilon},
  \qquad \varepsilon = 10^{-8},
  \label{eq:grpo-adv}
\end{equation}
and the policy maximizes a clipped surrogate against a KL penalty toward a
reference policy $\polref$. With $\rho_{i,t}$ the per-token importance ratio,
the implemented loss is
\begin{equation}
  \mathcal{L}^{\mathrm{GRPO}}(\theta)
  = -\frac{1}{\Ggrp L}\sum_{i=1}^{\Ggrp}\sum_{t=1}^{L}
      \min\!\bigl(\rho_{i,t} A_i,\;
        \mathrm{clip}(\rho_{i,t}, 1-\epsilon_c, 1+\epsilon_c) A_i \bigr)
    \;+\; \beta_{\mathrm{KL}}\, \widehat{\mathrm{KL}}_{i,t},
  \label{eq:grpo}
\end{equation}
with $\epsilon_c = 0.2$ and the standard $k_3$ estimator, written in terms of
$u_{i,t} = \log \polref(z_{i,t}) - \log \pol(z_{i,t})$,
\begin{equation}
  \widehat{\mathrm{KL}}_{i,t} \;=\; e^{u_{i,t}} - u_{i,t} - 1 \;\ge\; 0 .
  \label{eq:k3}
\end{equation}
Because prompts have fixed length $L$, the token mean in Eq.~\eqref{eq:grpo}
carries no length bias, so the length-normalization critique levelled at GRPO by
Dr.\ GRPO does not apply in this setting. An optional entropy bonus
$-\,\alpha_H H(\pol)$ may be added; it is off in the plain baseline.

\paragraph{The clip is provably inert here, and this matters.}
With a single update per rollout, the behaviour policy \emph{is} the current
policy, so $\rho_{i,t} \equiv 1$ identically --- both in value and, since $1$
lies strictly inside $[1-\epsilon_c, 1+\epsilon_c]$, in gradient. The two
branches of the $\min$ therefore coincide, and Eq.~\eqref{eq:grpo} reduces to
REINFORCE with a group-normalized advantage plus a KL penalty:
\begin{equation}
  \nabla_\theta \mathcal{L}^{\mathrm{GRPO}}
  = -\frac{1}{\Ggrp L}\sum_{i,t} A_i \nabla_\theta \log \pol(z_{i,t})
    \;+\; \beta_{\mathrm{KL}} \nabla_\theta \widehat{\mathrm{KL}} .
\end{equation}
The consequence is structural rather than cosmetic: \textbf{the KL term is the
only thing regularizing this objective}. Its unconstrained optimum is a
deterministic policy --- put all mass on the group's best prompt --- so with the
trust region inactive, GRPO's fixed point in this pipeline is entropy collapse.
Implementations that run GRPO with $\beta_{\mathrm{KL}} = 0$ (TRL's default;
DAPO; Dr.\ GRPO) can do so because they take multiple epochs per rollout, which
keeps the clip live. This pipeline does not, so we set
$\beta_{\mathrm{KL}} = 0.04$, the canonical DeepSeekMath value.

\paragraph{The reference anchor must be long-lived.}
A KL penalty only regularizes if the anchor is stable. If $\polref$ is
re-copied from the live policy every $k$ steps, then each refresh sets both the
penalty and its gradient to zero, and for small $k$ the anchor never accumulates
restoring force --- the policy drifts arbitrarily far from initialization while
staying within $\sim\!1\%$ of a reference that chases it. We therefore
\emph{pin} the anchor for the entire run ($k = 0$: no re-anchoring), matching
TRL's \texttt{sync\_ref\_model=False} default, verl's and OpenRLHF's frozen
reference, and DeepSeekMath's once-per-outer-iteration re-anchoring.
Section~\ref{sec:baseline-defense} reports what this fix is worth empirically.

\begin{algorithm}[t]
\caption{GRPO as implemented (one update per rollout, pinned KL anchor)}
\label{alg:grpo}
\begin{algorithmic}[1]
\STATE \textbf{input} policy $\pol$, group size $\Ggrp$, KL weight $\beta_{\mathrm{KL}}$, clip $\epsilon_c$
\STATE $\polref \leftarrow \pol$ \COMMENT{pinned once; never re-anchored}
\FOR{step $= 1$ \TO $T$}
  \STATE draw a batch of sources $x^{(1..B)}$ and $\lmb^{(b)} \sim \mathcal{U}[0,1]$
  \FOR{each source $b$}
    \STATE sample $z_1,\dots,z_\Ggrp \sim \pol(\cdot \mid \lmb^{(b)})$
    \STATE $r_i \leftarrow \hat{r}_{\lmb^{(b)}}(z_i, x^{(b)})$ \COMMENT{avg.\ of $N=50$ frozen-task-LM samples}
    \STATE $A_i \leftarrow (r_i - \mathrm{mean}(r))/(\mathrm{std}(r) + \varepsilon)$
  \ENDFOR
  \STATE $\mathcal{L} \leftarrow$ Eq.~\eqref{eq:grpo} averaged over all groups and tokens
  \STATE one gradient step on $\theta$ \COMMENT{$\Rightarrow \rho \equiv 1$, clip inert}
\ENDFOR
\end{algorithmic}
\end{algorithm}

\subsection{\rrebel{}: robust pairwise reward-difference regression}
\label{sec:rrebel}

\rrebel{} discards the advantage-and-clip machinery entirely and treats policy
improvement as a \emph{regression} problem. Define the $\beta$-scaled log-ratio
of a prompt against a reference policy,
\begin{equation}
  \logratio(z) \;=\; \beta\Bigl(\log \pol(z \mid \lmb) - \log \polref(z \mid \lmb)\Bigr).
  \label{eq:logratio}
\end{equation}
For every \emph{unordered pair} $(i,j)$, $i<j$, within a group, \rrebel{} asks
that the model's log-ratio difference match the observed reward difference, and
penalizes the mismatch with a discrepancy $d$:
\begin{equation}
  \mathcal{L}^{\rrmath}(\theta)
  \;=\; \frac{2}{\Ggrp(\Ggrp-1)} \sum_{i<j}
        d\Bigl(
          \underbrace{\logratio(z_i) - \logratio(z_j)}_{\text{model}}
          \;-\;
          \underbrace{(\tilde{r}_i - \tilde{r}_j)}_{\text{target}}
        \Bigr),
  \label{eq:rrebel}
\end{equation}
where $\tilde{r}$ is the preprocessed reward described below. Taking
$d(\cdot) = (\cdot)^2$ recovers REBEL, whose least-squares fixed point is the
tempered softmax
\begin{equation}
  \pi^\star(z) \;\propto\; \polref(z)\,\exp\!\bigl(r(z)/\beta\bigr).
  \label{eq:fixedpoint}
\end{equation}
This is the structural difference from GRPO, and the paper's central mechanistic
claim: \textbf{\rrebel{}'s optimum is a distribution, not a point mass.}
Exploration does not have to be bolted on by a KL penalty or an entropy bonus,
because a policy that collapses cannot satisfy Eq.~\eqref{eq:rrebel} --- a
deterministic policy makes every log-ratio difference infinite while the reward
differences stay finite and bounded. Since a $\lmb$-conditioned policy must
serve the entire curve, this matters more here than in single-objective RLHF: a
collapsed policy is not merely suboptimal, it is $\lmb$-independent and yields a
point instead of a frontier.

\paragraph{Robust discrepancies.}
The \emph{R} in \rrebel{} is the replacement of the squared loss by a robust
$d$. The motivation is that $\tilde{r}_i - \tilde{r}_j$ is a Monte Carlo
estimate: with $N=50$ stochastic task-model samples per prompt, a heavy right
tail in the sentiment classifier produces occasional large, spurious reward
differences, and under a squared loss a single such pair dominates the group's
gradient. We consider
\begin{equation}
  d(e) \in \Bigl\{\, e^2,\quad |e|,\quad
    \mathrm{Huber}_\delta(e),\quad
    \log\bigl(1 + \tfrac{e^2}{2c^2}\bigr) \,\Bigr\}
\end{equation}
--- squared ($\ell_2$, i.e.\ REBEL), absolute ($\ell_1$),
Huber with $\delta = 1$, and Cauchy with $c^2 = 1/2$ --- whose influence
functions are respectively unbounded, bounded, bounded, and redescending.

\paragraph{Reward preprocessing, and why it is scale-fair.}
Rewards are scaled by a constant and a $\lmb$-dependent factor, and then
optionally standardized within the group:
\begin{equation}
  \tilde{r}_i \;=\; \frac{\kappa\, g(\lmb)\, r_i}
                         {\mathrm{std}(\kappa g(\lmb) r_{1:\Ggrp}) + \varepsilon},
  \qquad \kappa = 0.1, \quad g(\lmb) = 2.2111 \cdot 10^{\lmb} - 2.1111 .
  \label{eq:rewardprep}
\end{equation}
The group standardization in the denominator \emph{exactly cancels} any positive
per-group scale, so for every \rrebel{} variant reported in this paper (all of
which use it) the constants $\kappa$ and $g(\lmb)$ have no effect on the
gradient. This is a deliberate fairness property rather than an accident: GRPO's
advantage normalization in Eq.~\eqref{eq:grpo-adv} cancels reward scale in the
same way, so neither family can benefit from a tuned reward scale that the other
lacks. Both families therefore regress or ascend on a scale-normalized target.

\paragraph{Reference policy.}
Unlike GRPO's pinned anchor, \rrebel{}'s $\polref$ is refreshed every
$150$ steps. This is not a hyperparameter choice but part of REBEL's
construction: Eq.~\eqref{eq:rrebel} regresses the log-ratio against reward
differences measured under the \emph{recent} policy, so the reference must lag
by a bounded amount. The asymmetry with GRPO is intentional and is exactly the
distinction that the earlier defect blurred (Section~\ref{sec:baseline-defense}).

\paragraph{Variants evaluated.}
We report three: $\ell_1$-std ($d = |\cdot|$, group standardization),
Huber-std ($d = \mathrm{Huber}_1$, group standardization), and
$\ell_1$-ent ($\ell_1$-std plus an entropy bonus $-\alpha_H H(\pol)$,
$\alpha_H = 0.01$). Plain $\ell_2$ (REBEL) was dropped after ranking fourth and
proving seed-unstable in an earlier campaign.

\begin{algorithm}[t]
\caption{\rrebel{} (robust pairwise reward-difference regression)}
\label{alg:rrebel}
\begin{algorithmic}[1]
\STATE \textbf{input} policy $\pol$, group size $\Ggrp$, scale $\beta$, discrepancy $d$, ref lag $K = 150$
\STATE $\polref \leftarrow \pol$
\FOR{step $= 1$ \TO $T$}
  \IF{step $\bmod\ K = 0$}
    \STATE $\polref \leftarrow \pol$ \COMMENT{bounded lag, required by the construction}
  \ENDIF
  \STATE draw sources $x^{(1..B)}$ and $\lmb^{(b)} \sim \mathcal{U}[0,1]$
  \FOR{each source $b$}
    \STATE sample $z_1,\dots,z_\Ggrp \sim \pol(\cdot \mid \lmb^{(b)})$
    \STATE $r_i \leftarrow \hat{r}_{\lmb^{(b)}}(z_i, x^{(b)})$; standardize to $\tilde{r}_i$ by Eq.~\eqref{eq:rewardprep}
    \STATE $\logratio(z_i) \leftarrow \beta(\log \pol(z_i) - \log \polref(z_i))$
  \ENDFOR
  \STATE $\mathcal{L} \leftarrow$ Eq.~\eqref{eq:rrebel} over all $\binom{\Ggrp}{2}$ within-group pairs
  \STATE one gradient step on $\theta$
\ENDFOR
\end{algorithmic}
\end{algorithm}
